% Original analytic teaching diagram, 2026-11-09.
\definecolor{VB}{HTML}{7188A3}\definecolor{VBF}{HTML}{E5EAF0}
\definecolor{VG}{HTML}{829C8B}\definecolor{VGF}{HTML}{E8EEE9}
\definecolor{VY}{HTML}{B4A16F}\definecolor{VYF}{HTML}{F1ECDD}
\definecolor{VR}{HTML}{A7787D}\definecolor{VRF}{HTML}{F0E5E6}
\definecolor{VP}{HTML}{9588A5}\definecolor{VPF}{HTML}{ECE8F0}
\definecolor{VI}{HTML}{3B4252}\definecolor{VM}{HTML}{6B7280}\definecolor{VL}{HTML}{4C4D4F}
\begin{tikzpicture}[x=1mm,y=1mm,>=Stealth,font=\figurefont\fontsize{8}{10}\selectfont,
lab/.style={text=VM,font=\figurefont\fontsize{7.5}{9}\selectfont},
pt/.style={circle,inner sep=0pt,minimum size=2mm},
vec/.style={->,line width=.95pt},
box/.style={draw=VB,fill=VBF,rounded corners=1.5pt,minimum height=7.5mm,align=center,inner sep=2pt}]
% Relations among the four kinds of results.
\node[box,minimum width=33mm] (w) at (0,24) {加权和全局最优\\（权重全正）};
\node[box,minimum width=33mm,fill=VRF,draw=VR] (p) at (0,0) {全局Pareto最优};
\node[box,minimum width=33mm] (s) at (72,0) {Pareto驻定};
\node[box,minimum width=33mm,fill=VYF,draw=VY] (c) at (72,24) {有限非支配候选集};
\draw[->,VL,line width=.75pt] (w)--node[lab,right]{无需凸性} (p);
\draw[->,VL,line width=.75pt] ([yshift=1.6mm]p.east)--node[lab,above]{可微、约束资格} ([yshift=1.6mm]s.west);
\draw[->,VL,line width=.75pt,dash pattern=on 2pt off 1.5pt] ([yshift=-1.6mm]s.west)--node[lab,below]{另需凸性且权重全正} ([yshift=-1.6mm]p.east);
\draw[->,VR,line width=.75pt,dash pattern=on 2pt off 1.5pt] (c.south west)--node[pos=.5,fill=white,text=VR,inner sep=1pt,font=\figurefont\fontsize{7.5}{9}\selectfont]{不能推出} (p.north east);
\end{tikzpicture}
